\section{Survival of a set of assignments}
\label{sec:replacement}

The deletion argument will leave a set of assignments that must pass
random tests of length \(k-1\).  We need to relate the probability that
those tests eliminate the set to its lifetime under clauses of length
\(k\).  This section proves the relation for every assignment set.
Replacing a shorter constraint by a block of clauses of the original
length has a qualitative predecessor in Friedgut's half-cube
argument~\cite[Lemma~5.7]{Friedgut}.

Fix integers \(k\ge3\) and \(u\ge k\), and put
\[
 r=k-1,\qquad \beta=\frac{k}{k-1}.
\]
A proper \(s\)-clause, for \(1\le s\le u\), chooses an \(s\)-element
subset of the \(u\) variables uniformly and then chooses independent
fair signs on that subset.  Clauses
are independent, so a whole clause may occur more than once.  A
collection of clauses \emph{kills} \(S\subseteq\{0,1\}^u\) when no
member of \(S\) satisfies every clause in the collection.  For
nonnegative integers \(d,g\), define
\[
 q_d(S)=\PP(\text{\(d\) proper \(r\)-clauses kill \(S\)}),
 \qquad
 R_g(S)=\PP(\text{\(g\) proper \(k\)-clauses kill \(S\)}).
\]
For a nonempty set, \(q_0(S)=R_0(S)=0\).  For the empty set these
probabilities are one, including when there are zero clauses.
Let \(H_k(S)\) be the first killing index in an infinite sequence of
proper \(k\)-clauses, counting the killing clause, and put
\[
 H_k(\varnothing)=0,\qquad \tau_k(S)=\EE H_k(S).
\]
These expectations are finite: for every integer \(m\ge0\),
\begin{equation}\label{eq:set-survival}
 \PP(H_k(S)>m)
 \le |S|(1-2^{-k})^m
 \le 2^u\exp(-2^{-k}m).
\end{equation}
Indeed, each fixed assignment survives each independent clause with
probability \(1-2^{-k}\), and the first inequality is a union bound.

\subsection{Replacing one shorter clause}

For a nonempty set \(S\), a coordinate is \emph{frozen} if every member
of \(S\) has the same value there.  If \(b\) coordinates are frozen,
the probability that one proper \(s\)-clause kills \(S\) is
\begin{equation}\label{eq:frozen-killing}
 h_s(S)=\frac{(b)_s}{2^s(u)_s},\qquad 1\le s\le u,
\end{equation}
where \((a)_s=a(a-1)\cdots(a-s+1)\), with value zero for a nonnegative
integer \(a<s\).  To kill \(S\), every literal of the clause must be
false for every assignment in \(S\).  Its variables must therefore
be frozen, and each sign must oppose their common values.  This
characterizes the killing clauses and proves the formula.

Write \(h=h_r(S)\).  Then \(0\le h\le2^{-r}\le1/4\).  If \(b\ge k\),
\begin{equation}\label{eq:frozen-power}
 h_k(S)=h\frac{b-r}{2(u-r)}\ge\frac1k h^\beta.
\end{equation}
In fact \(b-r\ge b/k\), \(u-r\le u\), and
\[
 h=\frac1{2^r}\prod_{j=0}^{r-1}\frac{b-j}{u-j}
 \le \left(\frac b{2u}\right)^r.
\]
Substituting these bounds into the equality proves
\eqref{eq:frozen-power}.  Its hypothesis matters.  For the cylinder
whose first \(r\) coordinates are fixed to zero,
\begin{equation}\label{eq:frozen-exception}
 h_k(S)=0,\qquad h_r(S)=\frac1{2^r\binom ur}>0.
\end{equation}
No single \(k\)-clause kills this set.  A sufficiently long block
can still do so through the joint effect of its clauses.

\begin{lemma}[One-clause replacement]\label{lem:one-replacement}
For an integer \(k\ge3\), let
\[
 K_k=\bigl(2^k(k-1)\bigr)^{k-1}.
\]
For every integer \(u\ge k\), every \(S\subseteq\{0,1\}^u\), and
every integer \(g\ge1\),
\begin{equation}\label{eq:one-replacement}
 R_g(S)\ge q_1(S)-K_k g^{-(k-1)}.
\end{equation}
\end{lemma}

\begin{proof}
For \(S=\varnothing\), both probabilities are one.  Suppose \(S\) is
nonempty and retain \(b,r,h\) above.  A block kills \(S\) if any
one of its clauses does.  Independence gives
\begin{equation}\label{eq:individual-block-bound}
 R_g(S)\ge1-(1-h_k(S))^g\ge1-e^{-g h_k(S)}.
\end{equation}

First assume \(b\ge k\).  If \((g/k)h^{1/r}\ge2\), then
\eqref{eq:frozen-power} gives \(g h_k(S)\ge2h\), and
\[
 R_g(S)\ge1-e^{-2h}\ge h.
\]
For the last inequality, \(1-e^{-2h}-h\) is zero at \(h=0\) and
has nonnegative derivative on \([0,1/4]\).  If instead
\((g/k)h^{1/r}<2\), then
\[
 h<(2k)^r g^{-r}\le K_k g^{-r},
\]
because \(2k\le2^k(k-1)\).  Nonnegativity of \(R_g(S)\) now suffices.

If \(b<r\), then \(h=0\), which also proves the inequality.  The only
remaining case is \(b=r\).  Formula~\eqref{eq:frozen-killing} gives
\[
 h=\frac1{2^r\binom ur}\le\left(\frac r{2u}\right)^r.
\]
When \(g<2^{k+1}u\), this is less than
\((2^k r/g)^r=K_k g^{-r}\).  When \(g\ge2^{k+1}u\), apply
\eqref{eq:set-survival} to the full cube:
\[
 R_g(S)\ge1-\exp((\log2-2)u)
 \ge1-e^{-u}>\frac14\ge h.
\]
Here \(u\ge k\ge3\).  This last bound uses the whole block and remains
valid even though \(h_k(S)=0\).
\end{proof}

\subsection{Several tests and the expected lifetime}

Replacing tests one at a time requires a bound for arbitrary residual
sets: conditioning on the other clauses need not leave a set with any
special description.  The sequential conditioning follows the replacement
argument of Carenini~\cite[Theorem~4.3]{Carenini}.  The uniform additive
block estimate in Lemma~\ref{lem:one-replacement} takes the place of the
cutoff-dependent comparison of single-clause killing probabilities in
\cite[Lemma~6.1]{Carenini}; the exceptional case \(b=k-1\) above
rules out a cutoff-free positive power lower bound for \(h_k\) in
terms of \(h_r\).

\begin{lemma}[Sequential replacement]\label{lem:sequential-replacement}
For integers \(k\ge3\) and \(u\ge k\), every
\(S\subseteq\{0,1\}^u\), and all integers \(d,g\ge1\),
\begin{equation}\label{eq:sequential-replacement}
 R_{dg}(S)\ge q_d(S)-dK_k g^{-(k-1)}.
\end{equation}
\end{lemma}

\begin{proof}
Choose independent proper \(r\)-clauses \(Q_1,\ldots,Q_d\) and
independent blocks \(G_1,\ldots,G_d\), each containing \(g\) proper
\(k\)-clauses; take all these choices mutually independent.  Let
\(p_j\) be the probability that
\(G_1,\ldots,G_j,Q_{j+1},\ldots,Q_d\) kill \(S\), for \(0\le j\le d\).

For the replacement at index \(j\), condition on
\(G_1,\ldots,G_{j-1},Q_{j+1},\ldots,Q_d\), and let \(S'\) be the
members of \(S\) satisfying those clauses.  The omitted \(Q_j\) and
fresh \(G_j\) remain independent with their original laws.  Their
conditional killing probabilities are \(q_1(S')\) and \(R_g(S')\).
Lemma~\ref{lem:one-replacement} applies also if \(S'=\varnothing\),
so \(p_j\ge p_{j-1}-K_k g^{-r}\) after averaging.  Summing from
\(j=1\) to \(d\), using \(p_0=q_d(S)\) and \(p_d=R_{dg}(S)\), proves
the claim.
\end{proof}

\begin{proposition}[Lifetime from shorter-clause killing]
\label{prop:lifetime}
For an integer \(k\ge3\), put
\[
 \beta=\frac{k}{k-1},\qquad
 K_k=\bigl(2^k(k-1)\bigr)^{k-1},\qquad
 C_k=4(2K_k)^{1/(k-1)}.
\]
For every integer \(u\ge k\), every \(S\subseteq\{0,1\}^u\), and
every integer \(d\ge1\),
\begin{equation}\label{eq:lifetime-product}
 \tau_k(S)q_d(S)^\beta\le C_k d^\beta.
\end{equation}
Consequently, for nonempty \(S\) and \(0<a\le q_d(S)\),
\begin{equation}\label{eq:lifetime-level}
 \tau_k(S)\le C_k d^\beta a^{-\beta}.
\end{equation}
There is no upper restriction on \(d\) and no positive lower cutoff
on \(a\).
\end{proposition}

\begin{proof}
For the empty set the lifetime is zero.  If \(S\ne\varnothing\) and
\(q_d(S)=0\), the left side of \eqref{eq:lifetime-product} is zero,
since the lifetime is finite.  Otherwise put \(q=q_d(S)>0\) and choose
\[
 g=\left\lceil\left(\frac{2K_kd}{q}\right)^{1/(k-1)}\right\rceil.
\]
Sequential replacement gives \(R_{dg}(S)\ge q/2\).  Split an infinite
independent sequence of \(k\)-clauses into blocks of length \(dg\).
The events that these blocks kill the \emph{original} set \(S\) are
independent, each of probability at least \(q/2\).  The first successful
block has expected index at most \(2/q\), and the cumulative sequence
kills \(S\) no later than the end of that block.  Hence
\[
 \tau_k(S)\le\frac{2dg}{q}.
\]
The quantity inside the ceiling is at least one, so \(g\) is at most
twice that quantity.  It follows that
\[
 \tau_k(S)
 \le4(2K_k)^{1/(k-1)}
       \left(\frac dq\right)^{1+1/(k-1)}
 =C_kd^\beta q^{-\beta}.
\]
This proves the product bound and then the stated level bound.
\end{proof}

These inequalities also hold conditionally for a set \(S\) measurable
with respect to an exposed history, provided the comparison clauses
remain independent of that history.  Apply the deterministic-set
inequality at each realization of \(S\).  In particular, sequential
replacement conditions only on other independent clause positions;
it does not condition the fresh clauses on a satisfiability event.
